\section{Le moteur de recherche contre celui qui exécute}
\label{sec:comparaison}

Deux implémentations indépendantes du même moteur coexistent~: celle de
recherche, en Python, et celle d'exécution, en MQL5, qui tourne dans le
simulateur du courtier. Les faire diverger est normal --- elles ne lisent pas
les mêmes données et ne remplissent pas les ordres de la même façon. Ce qui
importe est de savoir \emph{où} elles divergent et \emph{combien} cela coûte.

\subsection{Ce que la réconciliation avait conclu, et ce que le détail montre}

Le travail de réconciliation mené en juillet~2026 avait relevé
\textbf{35~trades de part et d'autre} sur la fenêtre commune et en avait
conclu que «~le signal et les transitions concordent~». L'examen transaction par
transaction confirme la première moitié de cette conclusion et corrige la
seconde.

\begin{figure}[H]
\centering
\includegraphics[width=\linewidth]{gold_trades_alignment}
\caption{Périodes durant lesquelles chaque moteur détenait une position. Le
fond rosé marque les jours où les deux moteurs ne sont pas dans le même état.}
\end{figure}

Au niveau des \textbf{fenêtres}, la concordance est excellente~: les deux
moteurs sont dans le même état \metric{90\,\%} des jours, pour une exposition
au marché identique de 75\,\%. Le signal est donc bien le même, et les périodes
où le moteur choisit de rester à l'écart --- 2022 notamment --- sont retrouvées
des deux côtés.

Au niveau des \textbf{transactions}, en revanche, l'appariement ne tient pas.
Sur les trente-cinq positions du simulateur~:

\begin{itemize}[leftmargin=1.4em, itemsep=0.3em]
\item \textbf{18} correspondent à une seule position côté recherche~;
\item \textbf{10} sont \emph{fractionnées} en deux ou trois positions côté
      recherche~;
\item \textbf{7} n'ont aucun équivalent, et symétriquement \textbf{8}
      positions de recherche n'existent pas côté exécution.
\end{itemize}

L'égalité des comptes est donc une coïncidence arithmétique, pas une
correspondance. Cela n'invalide pas la réconciliation --- son objet était de
vérifier que le signal se transporte, et il se transporte --- mais cela
interdit de comparer les trades un à un, et c'est pourquoi ce document compare
des états et des distributions.

\subsection{Deux profils de trade}

\begin{figure}[H]
\centering
\includegraphics[width=\linewidth]{gold_trades_engine_profiles}
\caption{Distribution des durées et des rendements de prix, moteur par moteur,
sur la même fenêtre.}
\end{figure}

Les deux distributions sont remarquablement semblables~: même durée médiane
(5~jours), même durée moyenne (34~jours), même durée maximale à un jour près, et
surtout un temps cumulé en position pratiquement identique --- \metric{1\,179}
jours-position pour l'exécution contre \metric{1\,178} pour la recherche. Les
deux moteurs passent donc exactement autant de temps investis.

La différence tient uniquement au \textbf{découpage} de ce temps. Là où le
moteur d'exécution tient une position d'un seul tenant, celui de recherche la
fractionne parfois en plusieurs allers-retours~: son score repasse brièvement
sous zéro, il sort, puis rentre à nouveau. Le cas le plus net est la position
ouverte le 24~novembre~2022 et tenue 89~jours à l'exécution, que la recherche
découpe en trois.

Ce détail a une conséquence qui dépasse la comptabilité, car le fractionnement
frappe précisément les positions longues --- celles qui produisent le résultat.
Les deux plus gros trades du moteur, tenus 214 et 199~jours, sont l'un et
l'autre fractionnés côté recherche. Or chaque aller-retour supplémentaire coûte
un écartement de prix sans rien apporter au signal.

\begin{insightbox}
On s'attend d'ordinaire à ce que le simulateur du courtier dégrade les
résultats de la recherche --- écartement, glissement, arrondis. Ici, sur ce
point précis, l'avantage va dans l'autre sens~: à temps d'exposition égal, le
moteur d'exécution paie \emph{moins} de rotations que celui de recherche. C'est
cohérent avec ce qui avait été observé lors du réglage des horizons, où les
horizons courts brillaient en recherche et s'effondraient à l'exécution.
\end{insightbox}

\subsection{Où passe l'argent}

\begin{figure}[H]
\centering
\includegraphics[width=\linewidth]{gold_trades_attribution}
\caption{Du résultat de prix au résultat net. Les trois postes sont mesurés sur
le journal des transactions et se recollent exactement au chiffre publié.}
\end{figure}

% Généré par scripts/build_gold_trades_figures.py — NE PAS ÉDITER À LA MAIN.
\begin{table}[H]
\centering
\begin{tabular}{@{}lrl@{}}
\toprule
Poste & USD & Source \\
\midrule
Résultat de prix (mesuré) & +44,926 & CSV des deals \\
Swap (mesuré) & -12,784 & 28.5\,\% du brut \\
Commission (mesurée) & +0 & nulle sur ce compte \\
\textbf{Résultat net} & \textbf{+32,141} & recolle au rapport MT5 \\
\bottomrule
\addlinespace[0.4em]
\multicolumn{3}{@{}p{0.92\linewidth}@{}}{\footnotesize Slippage estimé à 545\,\$ (2\,bps par côté sur 1.4\,M\$ de notionnel cumulé) : il est déjà contenu dans les prix d'exécution, ce n'est pas un poste additionnel. Le stop de sécurité a coupé 1 position pour -1,955\,\$. L'arrondi de lots au pas de 0,01 n'est pas chiffrable depuis les artefacts disponibles et reste \emph{non attribué}.}
\end{tabular}
\caption{Décomposition du résultat : ce qui est mesuré, ce qui est estimé}
\end{table}


Le résultat de prix s'élève à $+$44\,926\,\$. Le portage en retire
$-$12\,784\,\$, soit \metric{28,5\,\%}. La commission est nulle sur ce compte.
Il reste $+$32\,141\,\$, qui est exactement le chiffre du rapport
d'exécution --- l'attribution est complète, sans résidu.

\begin{figure}[H]
\centering
\includegraphics[width=\linewidth]{gold_trades_costs}
\caption{Résultat de prix et résultat net cumulés~: l'écart entre les deux
courbes est le portage. À droite, le portage payé par trade en fonction de sa
durée.}
\end{figure}

Le panneau de droite explicite le mécanisme~: le portage est un \emph{loyer de
temps}. Il croît avec la durée de détention, et frappe donc le plus lourdement
les positions qui rapportent le plus. Le trade de 214~jours paie 2\,874\,\$ de
portage, soit 11,7\,\% de son gain de prix~; celui de 199~jours en paie
2\,844\,\$, soit 18,9\,\% du sien. Rapporté à l'ensemble, le portage représente
28,5\,\% du résultat de prix --- davantage que sur les grands trades, parce que
les positions perdantes paient elles aussi leur loyer sans rien rapporter.

\begin{warningbox}
Ce poste n'est \textbf{pas modélisé} du côté recherche. Toute mesure de
performance issue du moteur de recherche surestime donc le résultat de
28\,\% environ sur ce moteur, avant même de considérer les autres écarts
d'exécution. C'est la raison pour laquelle le rapport d'investissement publie
les chiffres du simulateur et non ceux de la recherche.
\end{warningbox}

Deux postes plus fins complètent le tableau, et il faut être clair sur leur
statut. Le \textbf{glissement} est estimé à environ 545\,\$ sur 1,4~million de
dollars de notionnel cumulé, à partir du paramètre de 2~points de base par
côté~: cette somme est déjà contenue dans les prix d'exécution, ce n'est pas un
poste à retrancher en plus. L'\textbf{arrondi des volumes} au pas de 0,01~lot
n'est pas reconstructible à partir des artefacts disponibles~; il reste
\emph{non attribué}, et ce document ne lui invente pas de chiffre.
